% Compact optional insertion for the existing ALPHA-comparison appendix.
% Numerical methods and raw data belong in the reproducibility package.
% Existing bibliography keys ALPHA:2022rxj and ALPHA:2026vwc are retained.
\paragraph{Finite collection and transport.}
To make the extraction question quantitative, we compare the finite radiating
sheet enclosure with a source-driven finite Drude plasma of the same active
volume, wire filling ratio and surface-resistance law. The port model includes
localized Gaussian pickup columns, their electromagnetic backaction, and an
ideal coherent network feeding one receiver. Both reactive loading and
extraction strength are searched, and each loaded resonance is retuned to the
comparison mass. The number of pickup elements is not equated to the number
of noisy receivers. This is a reduced model of distributed extraction, not a
prediction for ALPHA's hardware.

For a $10\,\mathrm{cm}\times10\,\mathrm{cm}$ face and the stated halo,
the radiative-to-port scan-response ratios with four pickup columns are
approximately $0.55$, $1.94$, $14.2$ and $87.5$ at $0.5$, $1$, $2$ and
$4\,\mathrm{meV}$. Increasing the pickup count removes this sparse-receiver
advantage. Against the strongest tested profile, including an ideal continuous
receiver, the ratios become $0.42$, $0.46$, $0.52$ and $0.54$.
Thus the useful experimental question is whether external optics can deliver
the signal with sufficiently lower transmission loss. If their additional
power transmission is 0.8, the radiative design reaches parity when the port
chain loses approximately $2.9$, $2.6$, $2.4$ and $2.3$\,dB, respectively,
at equal output noise. These are measurable thresholds, not assigned ALPHA
losses.

For illustration, a single-mode rectangular copper output guide with aspect
ratio 2:1, operating frequency 1.7 times cutoff, and the common resistance law
reaches these losses after about 153, 44, 13 and 4\,cm. Guide attenuation grows
as $\nu^{5/3}$ under these assumptions. Alternative transport can improve the
port design. If the same magnet instead permits a port-coupled plasma to fill
30\,mm of depth while the radiative device uses one emitting face, the required
port-chain loss rises to about 8.2--18.8\,dB over these masses. A prototype
should therefore test delivered response and usable magnetic volume together.
The new calculation motivates radiative extraction in a specified engineering
regime; it establishes no universal conversion advantage.
